\documentclass[11pt,letterpaper]{article}
\usepackage[T1]{fontenc}
\usepackage{lmodern,amsmath,amssymb,amsthm,microtype}
\usepackage[margin=1in]{geometry}
\usepackage[hidelinks]{hyperref}
\newtheorem{theorem}{Theorem}
\title{A local-field obstruction to a unit-only conversion of Parts's scaffold}
\author{Aafi Khan Malek\\\small Independent Researcher, Ahmedabad, India\\
\small\href{mailto:aafimalek2032@gmail.com}{aafimalek2032@gmail.com}}
\date{November 2026}
\begin{document}
\maketitle
\begin{abstract}
The native coordinate field of Parts's \(\{1,2\}\)-distance scaffold
has a four-coloring of its entire unit-distance graph, with all
distance-two pairs monochromatic in the same coloring. The canonical
unit-chord gluing has the same property after adjoining a real square
root. A finite unit-distance gadget forcing unequal colors at distance
two must therefore leave these coordinate fields. This is a restricted
obstruction, not a new bound for the Euclidean plane.
\end{abstract}

\section{The native field and its four-coloring}
Let \(W\) consist of zero and the sixth roots of unity. In
\cite[\S3]{parts-two}, Parts uses
\[
 W\cup\alpha W\cup\alpha^{-1}W,\qquad
 \alpha=\exp\left(\tfrac i2\arccos(5/6)\right)
       =\frac{\sqrt{33}+\sqrt{-3}}6 ,
\]
where the square root of \(-3\) has positive imaginary part. The
Minkowski sums, the radius cut and the selected 199-point
half-scaffold all lie in \(K=\mathbb Q(\sqrt{33},\sqrt{-3})\).
The published obstruction uses both unit edges and distance-two
edges. We address the graph of unit edges alone.

\begin{theorem}
The unit-distance graph on all points of \(K\) has chromatic number
four. There is a four-coloring in which every distance-two pair is
monochromatic.
\end{theorem}
\begin{proof}
An odd integer congruent to one modulo eight is a square in
\(\mathbb Q_2\): starting from an odd root modulo eight, a root
modulo \(2^m\) can be adjusted by \(2^{m-1}\) to obtain a root
modulo \(2^{m+1}\). At each step the two candidate squares differ
by \(2^m\) modulo \(2^{m+1}\). The resulting compatible limit is a
square root. Apply this to 33, and embed \(K\) in
\(\mathbb Q_2(\omega)\), where \(\omega^2+\omega+1=0\), by mapping
\(\sqrt{-3}\) to \(1+2\omega\) and \(\sqrt{33}\) to a chosen
root in \(\mathbb Q_2\). The ring \(R=\mathbb Z_2[\omega]\) has
residue field \(\mathbb F_4\), since the polynomial is irreducible
modulo two. For an element \(a+b\omega\), its valuation is
\(\min(v_2(a),v_2(b))\): after division by the common power of
two, its residue is nonzero, and the product of two such residues
is nonzero. This gives a multiplicative valuation. The involution
\(\omega\mapsto-1-\omega\) preserves \(R\) and the valuation.
On the embedded \(K\) it agrees with ordinary complex conjugation.

If \(d\in K\) is a unit displacement then \(d\overline d=1\),
so \(v(d)=0\). If its length is two then \(v(d)=1\).
Write \(O\) for the inverse image of \(R\). In each additive coset
of \(O\) choose a representative and color the point by the
\(\mathbb F_4\)-residue of its difference from that representative.
Unit endpoints belong to the same coset and have distinct residues.
Distance-two endpoints belong to the same coset and have equal
residues. This is the standard valuation and coset method
(see, for example, Madore \cite[Prop.~3.2]{madore}), here applied
to a quadratic norm compatible with complex conjugation.

For sharpness, take two unit diamonds with shared endpoint zero,
the first with other endpoint \(u=\sqrt{-3}\) and middle vertices
\((u+1)/2,(u-1)/2\). Multiply the first diamond by
\(\rho_0=(5+\sqrt{-11})/6\in K\) to obtain the second, with other
endpoint \(v=\rho_0u\). The seven vertices are distinct:
the middle vertices have unit radius and angular separation
\(60^\circ\) within a diamond; the rotation is strictly between
zero and \(60^\circ\). The other endpoints have radius
\(\sqrt3\) and \(|u-v|^2=3(2-2\cdot5/6)=1\).
A proper three-coloring of each diamond makes its endpoints
the same color, contradicting this additional unit edge. This
is the ordinary Moser spindle.
\end{proof}

\section{Rotated copies and the missing gadget}
The two-distance construction in \cite[\S3]{parts-two} supplies a
pair at distance \(5/\sqrt3\) whose equality is forced under its
\emph{two-distance} five-color constraints. Two copies around a
shared terminal make a unit chord by the rotation
\[
 \gamma=\frac{47+\sqrt{-291}}{50},
 \qquad \operatorname{Re}\gamma=1-\frac{1}{2(25/3)}.
\]
The resulting points lie in
\(K'=\mathbb Q(\sqrt{33},\sqrt{97},\sqrt{-3})\). Since 97 is also
one modulo eight, the same embedding and valuation proof gives a
four-coloring of the unit edges on all of \(K'\), with every
distance-two pair monochromatic.

The separate unit-distance minimization work of Parts
\cite[\S\S2,4]{parts-unit} uses base coordinates in \(K\) but
joins base graphs using a different rotation,
\((7+\sqrt{-15})/8\). This multiplier leaves \(K\). The square
root of \(-15\) lies in \(\mathbb Q_2\), since \(-15\equiv1\pmod8\).
Its two choices are interchanged by complex conjugation, so that
conjugation cannot be the same field automorphism fixing
\(\mathbb Q_2\) under this enlarged embedding. The two valuations
of the multiplier are \(\pm2\): the numerators sum to 14 and
multiply to 64, so their valuations are 1 and 5. There is no
conflict with the unit-only five-chromatic joined graphs.

To convert the \(\{1,2\}\)-distance obstruction, one sufficient
ingredient is a finite graph of unit edges with exact terminals
\(u,v\) at distance two which must have different colors in every
proper five-coloring. For each distance-two edge \(a,b\) of the
scaffold, place a copy of the gadget by the isometry
\[
 z\longmapsto a+(b-a)(z-u)/(v-u).
\]
A proper five-coloring would give a proper two-distance
five-coloring of the original scaffold, which is impossible.
A vertex-minimal non-five-colorable subgraph of the resulting
finite graph would have chromatic number exactly six. The
required gadget has not been found here. The theorem rules out
any realization of it contained in \(K\) or \(K'\), up to isometry.

The proof applies at a conjugation-stable nonarchimedean place of
a complex number field whenever its residue field has at most
four elements. Ramification does not change the argument. At
residue characteristic two, the same coloring is monochromatic
on distance-two pairs. For the unramified 2-adic realization
above, a rational squared length of odd valuation cannot occur
at all, since \(v(d\overline d)=2v(d)\) is even.

\paragraph{A field-mixing test, not run.}
Adjoining \(\sqrt{-7}\) and using \(\beta=(3+\sqrt{-7})/4\) breaks
the specific 2-adic obstruction: the two possible valuations
of this complex unit are \(1,-1\). But
\(\mathbb Q(\sqrt{-3},\sqrt{-7})\) alone has a proper three-coloring.
At the unique 3-adic place, the residue field is
\(\mathbb F_9=\mathbb F_3[i]\), \(i^2=-1\), and conjugation sends
\(i\) to \(-i\). A unit displacement's residue \(a+bi\) satisfies
\(a^2+b^2=1\), so exactly one coefficient is zero and the color
difference \(a+b\) is nonzero. As before color each integral
coset separately. For a non-unit terminal pair, recoloring both
terminals with a fresh fourth color shows this smaller field
cannot supply the gadget. A proposed mixed fixture
\((W+W)+\alpha W+\beta W\), with at most 931 sites and terminals
\(0,2\), is not covered by those two coloring obstructions.
Its proper five-colorability subject to equal terminal colors
has not been tested. See the accompanying notes for the exact
structural proof and experiment scope.

\paragraph{Disclosure.} This written research note was developed
with AI assistance. It is not a Lean verification or a new
whole-plane chromatic bound.

\begin{thebibliography}{9}
\bibitem{parts-two} J. Parts, A small 6-chromatic two-distance graph
in the plane, \href{https://arxiv.org/abs/2010.12656v2}{arXiv:2010.12656v2}.
\bibitem{parts-unit} J. Parts, Graph minimization, focusing on the example
of 5-chromatic unit-distance graphs in the plane,
\href{https://arxiv.org/abs/2010.12665v2}{arXiv:2010.12665v2}.
\bibitem{madore} D. A. Madore, The Hadwiger--Nelson problem over certain
fields, \href{https://arxiv.org/abs/1509.07023}{arXiv:1509.07023}.
\end{thebibliography}
\end{document}
